%% Migrated from v1.1 Part IX, lines 12249--12969.  Statements only;
%% proved and sketched material is carried by annex/p02/higher-order.tex.
%% Phase-2 repairs retained in substance: R-031 (the injectivity induction in
%% full PA2 categoricity) and R-033 (real categoricity remains cited rather
%% than receiving an unsupported proof).  Source exercises are retired: their
%% substantive warnings and calculations are incorporated below.

\chapter{Second- and Higher-Order Logic}
\label{sec:p02-second-and-higher-order-logic}

Second-order syntax permits quantification over relations on the object
domain.  That syntax supports sharply different semantic regimes: full
semantics quantifies over every relation of the required arity, whereas
Henkin semantics quantifies over selected relation domains.  Simple type
theory repeats the distinction at function types.  Every consequence and
completeness assertion below therefore names both its syntax and its semantic
class \cite{Henkin1950,Manzano1996,BarwiseFeferman1985}.

\section{Relational second-order syntax}
\label{sec:p02-sol-syntax}

\begin{hbdefinition}{def:p02-sol-language}{Relational second-order syntax}
A \emph{relational second-order signature} \(\Lang_2\) contains the object
constants, object function symbols, and predicate constants of a first-order
signature.  For every integer \(n\geq1\), its logical variables include a
countably infinite stock \(\mathsf{RelVar}_n\) of \(n\)-ary relation
variables.  Object terms are the ordinary first-order terms.  The atomic
formulas are
\[
  t=u,\qquad P^n(t_1,\ldots,t_n),\qquad X^n(t_1,\ldots,t_n),
\]
where \(P^n\) is a predicate constant and
\(X^n\in\mathsf{RelVar}_n\).  Formulas are closed under the classical
connectives, object quantifiers, and
\(\forall X^n,\exists X^n\).  There are no function variables, relation
terms other than variables and predicate constants, or relation equality.
A \emph{second-order sentence} has no free object-variable or
relation-variable occurrence.
\index{second-order logic!relational language}
\index{relation variable}
\end{hbdefinition}

The restriction to relation variables is load-bearing: a formula may replace
a free occurrence of \(X^n\) only through a separate, capture-avoiding
relation-formula substitution operation.  The displayed calculus instead
instantiates relation quantifiers by relation variables or predicate
constants; formula-defined relations enter through comprehension.

\section{Full semantics and its exact metatheorems}
\label{sec:p02-sol-full-semantics}

\begin{hbdefinition}{def:p02-sol-full-semantics}{Full relational second-order semantics}
Let \(\mathcal M\) have nonempty object domain \(M\).  A full assignment
maps each object variable to an element of \(M\) and every \(n\)-ary
relation variable to an arbitrary member of \(\mathcal P(M^n)\).  Besides
the first-order clauses,
\[
\begin{aligned}
 \mathcal M,s\entails X^n(t_1,\ldots,t_n)
 &\quad\Longleftrightarrow\quad
 (\lbrack\!\lbrack t_1\rbrack\!\rbrack_s^{\mathcal M},\ldots,
  \lbrack\!\lbrack t_n\rbrack\!\rbrack_s^{\mathcal M})\in s(X^n),\\
 \mathcal M,s\entails\forall X^n\,\varphi
 &\quad\Longleftrightarrow\quad
 \mathcal M,s[X^n\mapsto R]\entails\varphi
 \text{ for every }R\subseteq M^n.
\end{aligned}
\]
The existential clause is dual.  Write
\(\Gamma\entails_{\mathrm{full}}\varphi\) when every such structure and
full assignment satisfying \(\Gamma\) satisfies \(\varphi\).  The word
\emph{full} specifies the relation-quantifier range; it does not say that
\(M\) is a set-theoretic universe.
\index{full second-order semantics}
\index{semantic consequence!full second-order}
\end{hbdefinition}

\begin{hbtheorem}{thm:p02-sol-full-comprehension}{proved=annex:p2.hol.1}{Full comprehension validity}
Let \(n\geq1\), let \(x_1,\ldots,x_n\) be pairwise distinct, and let
\(X^n\) not occur free in
\(\chi(x_1,\ldots,x_n,\bar p,\bar Q)\).  Under full semantics the
universal closure over the displayed object and relation parameters of
\[
 \exists X^n\,\forall x_1\cdots\forall x_n
 \bigl(X^n(x_1,\ldots,x_n)\leftrightarrow
       \chi(x_1,\ldots,x_n,\bar p,\bar Q)\bigr)
\]
is valid.  The formula \(\chi\) may contain relation quantifiers, so this
is impredicative comprehension.
\end{hbtheorem}

\begin{hbdefinition}{def:p02-pa2}{Full second-order Peano axioms}
In the signature \(\{0,S\}\), \(\mathsf{PA}_2\) comprises
\[
 \forall x\,Sx\ne0,
 \qquad
 \forall x\,\forall y\,(Sx=Sy\to x=y),
\]
and the single second-order induction sentence
\[
 \forall X\bigl[(X(0)\land
   \forall x(X(x)\to X(Sx)))\to\forall x\,X(x)\bigr].
\]
When \(+\) and \(\times\) are present, add the four first-order recursion
axioms
\[
 \forall x\,(x+0=x),\qquad
 \forall x\,\forall y\,(x+S(y)=S(x+y)),
\]
\[
 \forall x\,(x\times0=0),\qquad
 \forall x\,\forall y\,(x\times S(y)=(x\times y)+x).
\]
This is still a finite axiom set; its categoricity assertion below is
specifically full-semantic.
\index{second-order arithmetic}
\index{categoricity!second-order arithmetic}
\end{hbdefinition}

\begin{hbtheorem}{thm:p02-pa2-categoricity}{proved=annex:p2.hol.2}{Full \(\mathsf{PA}_2\) categoricity}
Every \(\{0,S\}\)-structure satisfying \(\mathsf{PA}_2\) under full
semantics is isomorphic to \((\N,0,n\mapsto n+1)\).  Consequently any two
such structures are isomorphic.  With the recursive axioms for addition and
multiplication, the isomorphism preserves those operations as well.
\end{hbtheorem}

\begin{hbtheorem}{thm:p02-sol-full-noncompact}{proved=annex:p2.hol.3}{Full-semantic noncompactness}
Full relational second-order semantic consequence is not compact: there is
a set \(T\) of second-order sentences every finite subset of which has a full
model while \(T\) itself has no full model.
\end{hbtheorem}

\begin{hbtheorem}{thm:p02-sol-full-no-effective-completeness}{sketched=annex:p2.hol.4}{No effective completeness for full validity}
Suppose an effective calculus means one whose finite derivations are
computably checkable, hence whose theorem set is computably enumerable.
There is no such calculus \(\mathsf D\) satisfying, for every relational
second-order sentence \(\sigma\),
\[
 \mathsf D\derives\sigma
 \quad\Longleftrightarrow\quad
 \entails_{\mathrm{full}}\sigma.
\]
Equivalently, the full second-order validities are not computably enumerable.
\end{hbtheorem}

The last result is stronger than noncompactness and retains the source's
effectiveness qualifier.  It uses the effective reduction from true
first-order arithmetic to full validity and the non-enumerability of true
arithmetic; the omitted coding details are controlled by
\cite{BarwiseFeferman1985,Shapiro1991}.

\begin{hbtheorem}{thm:p02-sol-real-categoricity}{cited=Shapiro1991}{Full-semantic categoricity of \(\mathbb R\)}
In \(\{0,1,+,\times,<\}\), take the first-order ordered-field axioms and
\[
\begin{aligned}
 \forall X\Bigl[&\bigl(\exists x\,X(x)\land
   \exists y\,\forall x(X(x)\to x\leq y)\bigr)\\
 &{}\to \exists z\Bigl(\forall x(X(x)\to x\leq z)\\
 &\hspace{25mm}{}\land
   \forall y((\forall x(X(x)\to x\leq y))\to z\leq y)\Bigr)\Bigr].
\end{aligned}
\]
where \(u\leq v\) abbreviates \(u<v\lor u=v\) in this statement.  Under
full semantics, every model of these axioms is isomorphic to the ordered field
\(\mathbb R\); hence every such model is uncountable.
\end{hbtheorem}

The uncountability conclusion uses the ordinary set-theoretic theorem for
\(\mathbb R\) \cite{Jech2003}; the categoricity proof remains citation-only,
as required by the source repair R-033.  The cited proof sends each element of
one model to the element of the other determined by the same cut of their
common rational subfield and verifies that this map is an ordered-field
isomorphism \cite{Shapiro1991}.

\begin{hbcorollary}{thm:p02-sol-full-ls-failures}{proved=annex:p2.hol.6}{Failure of both L\"owenheim--Skolem directions under full semantics}
Full relational second-order semantics has neither the blanket upward nor the
blanket downward L\"owenheim--Skolem property.  The finite theory
\(\mathsf{PA}_2\) has an infinite full model but no uncountable full model;
the finite ordered-field plus least-upper-bound theory of
Theorem~\ref{thm:p02-sol-real-categoricity} has an infinite full model but no
countable full model.
\index{Lowenheim-Skolem theorem@L\"owenheim--Skolem theorem!failure in full second-order logic}
\end{hbcorollary}

\section[Relational Henkin semantics]{Relational Henkin semantics and the matched calculus}
\label{sec:p02-sol-henkin-semantics}

\begin{hbdefinition}{def:p02-sol-henkin-structure}{Relational Henkin structure}
A \emph{relational Henkin structure} is
\[
 \mathcal H=(M,(\mathcal R_n)_{n\geq1},\ldots),
\]
where \(M\ne\varnothing\), every
\(\varnothing\ne\mathcal R_n\subseteq\mathcal P(M^n)\), and every
\(n\)-ary predicate constant is interpreted by a member of
\(\mathcal R_n\).  An \(n\)-ary relation variable ranges over
\(\mathcal R_n\).  Satisfaction otherwise uses the clauses of
Definition~\ref{def:p02-sol-full-semantics}, with that restricted range.
A full structure is the special case
\(\mathcal R_n=\mathcal P(M^n)\) for every \(n\).
Write \(\Gamma\entails_{\mathrm H}\varphi\) when every relational Henkin
structure under every matching assignment satisfying \(\Gamma\) satisfies
\(\varphi\).  This unadorned subscript ranges over arbitrary selected
relation domains; it does not silently impose comprehension or choice.
\index{Henkin semantics!relational second-order}
\index{Henkin structure!relational}
\index{semantic consequence!Henkin second-order}
\end{hbdefinition}

Merely selecting nonempty relation domains is too weak for a calculus that
contains comprehension or choice.  The next definition names the exact model
class needed by the displayed sheet; the adjective ``matched'' records this
syntactic--semantic equality of schema instances and has no broader meaning.

\begin{hbdefinition}{def:p02-sol-matched-henkin}{Matched Henkin models}
A relational Henkin structure is a \emph{matched
comprehension--\allowbreak choice Henkin model}, abbreviated an
\(\mathsf{H}_{\mathsf{CC}}\)-model, iff it satisfies every instance of both
schema families below, with exactly the eligibility conditions displayed on
\hyperref[sheet:sol]{\texttt{sheet:sol}}:
\[
\begin{aligned}
\mathsf{Comp}_{\chi}^{n}&\defeq
 \operatorname{ucl}\!\left(\exists X^n\,\forall\bar x
   (X^n(\bar x)\leftrightarrow\chi(\bar x,\bar p,\bar Q))\right),\\
\mathsf{Choice}_{\chi}^{n}&\defeq
 \operatorname{ucl}\!\Bigl(\Bigl(\forall\bar x\,\exists y\,
 \chi(\bar x,y,\bar p,\bar Q)\Bigr)\to{}\\
&\hspace{18mm}\exists F^{n+1}\,\forall\bar x\,\exists y\,
 \Bigl(F^{n+1}(\bar x,y)\land\chi(\bar x,y,\bar p,\bar Q)\\
&\hspace{38mm}\land
 \forall z(F^{n+1}(\bar x,z)\to z=y)\Bigr)\Bigr).
\end{aligned}
\]
Here \(n\geq1\); in comprehension \(X^n\) is not free in \(\chi\);
in choice \(F^{n+1}\) is fresh and not free in \(\chi\); and
\(\operatorname{ucl}\) universally closes exactly the permitted object and
relation parameters \(\bar p,\bar Q\).  Comprehension may be impredicative.
Choice asserts an admissible single-valued selecting graph and adds no
function variable.  It is not redundant with comprehension: comprehension
collects a relation already defined by a formula, while choice selects one
witness in every nonempty fibre, a step not available in a choice-free
metatheory merely from definability.
\index{comprehension!second-order}
\index{choice!relational graph schema}
\end{hbdefinition}

Under full semantics, comprehension is automatic.  The choice schema is
full-valid only relative to enough metatheoretic choice to select from every
displayed family of nonempty fibres.  In a matched Henkin model, both schema
families hold by definition; this is a restriction on the model class, not a
consequence of \(\mathcal R_n\subseteq\mathcal P(M^n)\).

\begin{hbdefinition}{def:p02-sol-henkin-consequence}{Matched Henkin consequence}
Write \(\Gamma\entails_{\mathrm H}^{\mathsf{CC}}\varphi\) iff every
\(\mathsf H_{\mathsf{CC}}\)-model under every matching object-and-relation
assignment that satisfies \(\Gamma\) also satisfies \(\varphi\).  The
superscript records the comprehension--choice closure and is necessary to
distinguish this consequence relation from consequence over arbitrary
relational Henkin structures.
\index{semantic consequence!Henkin second-order!matched schemas}
\end{hbdefinition}

The calculus registry displays the exact four relation-quantifier rules and
the two schema families on
\hyperref[sheet:sol]{\texttt{sheet:sol}}.  It inherits the finite-multiset
sequent conventions and object-equality rules of
\hyperref[sheet:lk-eq]{\texttt{sheet:lk-eq}}; rules are not duplicated here.

\begin{hbdefinition}{def:p02-sol-incidence-translation}{Incidence translation}
For the exact language of Definition~\ref{def:p02-sol-language}, form a
many-sorted first-order language with object sort \(D\), relation sort
\(S_n\) for each \(n\geq1\), incidence predicate
\(E_n:S_n\times D^n\), and a relation-sort constant \(\widehat P:S_n\) for
each predicate constant \(P^n\).  Define the metalinguistic translation
\(\varphi\mapsto\varphi^\sharp\) recursively by
\[
\begin{aligned}
 (X^n(\bar t))^\sharp&\defeq E_n(X^n,\bar t),
 & (P^n(\bar t))^\sharp&\defeq E_n(\widehat P,\bar t),\\
 (\forall x\,\psi)^\sharp&\defeq\forall x{:}D\,\psi^\sharp,
 & (\exists x\,\psi)^\sharp&\defeq\exists x{:}D\,\psi^\sharp,\\
 (\forall X^n\,\psi)^\sharp&\defeq\forall X^n{:}S_n\,\psi^\sharp,
 & (\exists X^n\,\psi)^\sharp&\defeq\exists X^n{:}S_n\,\psi^\sharp.
\end{aligned}
\]
It leaves object terms and object equality unchanged and commutes with every
classical connective.  This is a recursively eliminable metalinguistic map,
not an object-language connective.  It is nonredundant with erasing the
second-order binders: relation application becomes incidence, predicate
constants become relation-sort objects, and the two quantifier ranges become
different sorts.  These distinctions are used in the correspondence and
adequacy theorems below.
For a formula set \(\Gamma\), write
\(\Gamma^\sharp\defeq\{\psi^\sharp:\psi\in\Gamma\}\); this is the
pointwise image under the same map, not an additional closure operation.
\index{incidence translation}
\index{second-order logic!many-sorted translation}
\end{hbdefinition}

Given \(\mathcal H\), define the strict many-sorted structure
\(\mathcal H^\sharp\) by tagging every sort:
\[
 D^{\mathcal H^\sharp}=\{D\}\times M,
 \qquad S_n^{\mathcal H^\sharp}=\{S_n\}\times\mathcal R_n,
\]
\[
\begin{aligned}
 E_n^{\mathcal H^\sharp}
   \bigl((S_n,R),(D,a_1),\ldots,(D,a_n)\bigr)
 &\Longleftrightarrow(a_1,\ldots,a_n)\in R,\\
 \widehat P^{\mathcal H^\sharp}
 &= (S_n,P^{\mathcal H}).
\end{aligned}
\]
Object constants and functions are transported along the object tag:
\(c^{\mathcal H^\sharp}=(D,c^{\mathcal H})\) and
\(f^{\mathcal H^\sharp}((D,\bar a))=(D,f^{\mathcal H}(\bar a))\).
For a matching assignment \(s\), put \(s^\sharp(x)=(D,s(x))\) and
\(s^\sharp(X^n)=(S_n,s(X^n))\).  Tagging is necessary because the sets
\(M,\mathcal R_1,\mathcal R_2,\ldots\) may overlap; it makes the incidence
structure satisfy the strict disjoint-sort convention used below.

\begin{hbtheorem}{thm:p02-sol-incidence-satisfaction}{proved=annex:p2.hol.7}{Incidence translation preserves satisfaction}
For every relational Henkin structure \(\mathcal H\), every second-order
formula \(\varphi\), and every matching assignment \(s\),
\[
 \mathcal H,s\entails\varphi
 \quad\Longleftrightarrow\quad
 \mathcal H^\sharp,s^\sharp\entails\varphi^\sharp.
\]
Moreover, \(\mathcal H\) is an \(\mathsf H_{\mathsf{CC}}\)-model iff
\(\mathcal H^\sharp\) satisfies the many-sorted translations of exactly all
eligible \(\mathsf{Comp}_{\chi}^{n}\) and
\(\mathsf{Choice}_{\chi}^{n}\) instances.
\end{hbtheorem}

\begin{hbdefinition}{def:p02-ms-variable-guard}{Total sorted-assignment guard}
For a strict many-sorted signature \(\Lang_{\mathrm{MS}}\), write
\(Q_{\mathsf s}\) for the unary predicate of its flat one-sorted translation
that marks sort \(\mathsf s\).  Because assignments are total, guarding only
variables free in the formulas would leave the values of every other sorted
variable to an unlicensed choice across the nonempty sort domains.  Define
the metalinguistic premise set
\[
 \mathsf{VarGuard}_{\mathrm{MS}}(\Lang_{\mathrm{MS}})
 \defeq\{Q_{\mathsf s}(v):\mathsf s\in\mathsf{Sort},
       \ v\text{ is a variable of sort }\mathsf s\}.
\]
This is an assignment guard, not an object-language operator or schema.  It
differs from the bound-variable guards inserted by
\(\operatorname{tr}_{\mathrm{MS}}\) and from
\(\mathsf{Ax}_{\mathrm{sort}}\), which constrains structures rather than
assignments.  The translated calculus uses the set immediately below, and
the general reduction theorem reuses it in
Section~\ref{sec:p02-many-sorted}.
\index{assignment guard!many-sorted}
\end{hbdefinition}

\begin{hbdefinition}{def:p02-sol-derived-calculus}{Translation-defined SOL calculus}
Let \(T_{\mathsf{CC}}^\sharp\) be the many-sorted first-order theory
consisting of the incidence translations of exactly the eligible
comprehension and choice instances above.  For a many-sorted formula set
\(\Delta\) and formula \(\psi\), define
\[
 \Delta\derives_{\mathsf{MSFOL}}\psi
 \quad\Longleftrightarrow\quad
 \mathsf{Ax}_{\mathrm{sort}}\cup
 \mathsf{VarGuard}_{\mathrm{MS}}(\Lang_{\mathrm{MS}})\cup
 \operatorname{tr}_{\mathrm{MS}}[\Delta]
 \derives_{\mathsf{ND}_{\mathsf{FOL}}}
 \operatorname{tr}_{\mathrm{MS}}(\psi),
\]
using Definitions~\ref{def:p02-many-sorted-translation} and the registered
first-order natural-deduction calculus.  This is again an eliminable
metalinguistic abbreviation; Theorem~\ref{thm:p02-many-sorted-translation},
its formula induction, and classical first-order adequacy make it sound and
strongly complete under the same set-theoretic hypotheses.  Now define
\[
 \Gamma\derives_{\mathsf{SOL}_{\mathrm H}^{\mathsf{CC}}}\varphi
 \quad\Longleftrightarrow\quad
 T_{\mathsf{CC}}^\sharp\cup\Gamma^\sharp
 \derives_{\mathsf{MSFOL}}\varphi^\sharp.
\]
The left side is thus an eliminable metalinguistic abbreviation for the right
side, not a newly postulated native derivation relation.  This definition is
needed because Henkin's typed completeness theorem does not cover the new
relational rule surface; it is used immediately to state an exact adequacy
theorem.  The formulas displayed on \texttt{sheet:sol} are admissible
projections of this definition, rather than an independently complete finite
list of native generators.

The advertised two-sided sheet has the following exact derived meaning.  For
finite-multiset contexts \(\Sigma,\Pi\), let \(\Sigma^\sharp\) be the
translated support of \(\Sigma\), and let \(\bigvee\Pi^\sharp\) be the
disjunction of any listing of \(\Pi^\sharp\), with
\(\bigvee\varnothing\defeq\bottom\).  Declare
\[
 \Sigma\Rightarrow\Pi\text{ derivable on \texttt{sheet:sol}}
 \quad\Longleftrightarrow\quad
 T_{\mathsf{CC}}^\sharp\cup\Sigma^\sharp
 \derives_{\mathsf{MSFOL}}\bigvee\Pi^\sharp.
\]
Classical first-order derivability makes this independent of listing,
bracketing, and repeated occurrences.  Thus it is an eliminable sequent
abbreviation, not a second calculus.  Its assignment guard includes every
sorted object and relation variable, hence in particular every variable free
in either context.
\index{second-order calculus!translation-defined}
\end{hbdefinition}

\begin{hbtheorem}{thm:p02-sol-derived-adequacy}{proved=annex:p2.hol.8}{Strong adequacy of the matched relational calculus}
Assume the Axiom of Choice in the ambient set theory.  For every set
\(\Gamma\) of relational second-order formulas and every formula \(\varphi\),
\[
 \Gamma\derives_{\mathsf{SOL}_{\mathrm H}^{\mathsf{CC}}}\varphi
 \quad\Longleftrightarrow\quad
 \Gamma\entails_{\mathrm H}^{\mathsf{CC}}\varphi.
\]
All six additions on \hyperref[sheet:sol]{\texttt{sheet:sol}} are admissible
for this derivability relation and for its two-sided projection just defined.
The conclusion is
strong completeness for the exact matched
selected-domain semantics; it is not completeness for full validity.
\end{hbtheorem}

\begin{proofboundary}{def:p02-sol-native-boundary}{What has and has not been proved about the displayed sheet}
The theorem defines derivability by the exact many-sorted reduction and then
proves the displayed rules admissible.  It does not assert that the inherited
first-order rules, four relation-quantifier rules, and two schema displays,
taken as an independently generated native calculus, are complete.  Such a
claim would additionally require a converse deductive-correspondence theorem.
Nor does \cite{Henkin1950} supply that missing theorem: it proves soundness and
completeness for Henkin's own typed system and its general models.
\end{proofboundary}

\begin{hbexample}{def:p02-sol-reachability}{Reachability under two semantics}
For a directed graph with edge relation \(E\), set
\[
 \operatorname{Reach}(a,b)\defeq
 \forall X\bigl[(X(a)\land
   \forall u\,\forall v((X(u)\land E(u,v))\to X(v)))\to X(b)\bigr].
\]
Under full semantics this holds exactly when a finite directed path leads
from \(a\) to \(b\): the converse direction instantiates \(X\) by the actual
set of vertices reachable from \(a\).  In an arbitrary Henkin structure that
set need not be admissible, so the converse can fail.
\end{hbexample}

\begin{hbwarning}{def:p02-sol-semantic-parameter}{Never omit the second-order semantic parameter}
The expressions \(\Gamma\entails_{\mathrm{full}}\varphi\) and
\(\Gamma\entails_{\mathrm H}^{\mathsf{CC}}\varphi\) name different
consequence relations.  Combining full categoricity with Henkin completeness
into the sentence ``second-order logic is complete and categorical'' is
incorrect: the two claims quantify over different model classes.
\end{hbwarning}

\section{Full versus Henkin second-order semantics}
\label{sec:p02-sol-comparison}

\begin{landscape}
\begingroup\scriptsize\setlength{\tabcolsep}{4pt}
\begin{longtable}{@{}p{1.35in}p{2.95in}p{2.95in}@{}}
\caption{Full versus matched Henkin relational semantics (apparatus A-039).}
\label{tab:p02-full-vs-henkin-sol}\\
\toprule
Property & Full relational semantics & Matched Henkin relational semantics \\ \midrule
\endfirsthead
\toprule
Property & Full relational semantics & Matched Henkin relational semantics \\ \midrule
\endhead
Relation-quantifier range & Every member of \(\mathcal P(M^n)\). & A selected nonempty \(\mathcal R_n\subseteq\mathcal P(M^n)\). \\
\addlinespace
Comprehension & Automatically valid for every eligible impredicative instance. & Must be imposed as satisfaction of exactly the eligible instances, equivalently the matching definability closure. \\
\addlinespace
Choice & Valid only relative to adequate metatheoretic choice. & Must be imposed by the eligible graph-choice instances or equivalent selector closure. \\
\addlinespace
Completeness & No sound effective recursively axiomatized calculus is complete for full validity. & Under ambient AC, strong completeness holds for the exact translation-defined calculus paired with the exact schema-closed model class, Theorem~\ref{thm:p02-sol-derived-adequacy}. \\
\addlinespace
Compactness & Fails, Theorem~\ref{thm:p02-sol-full-noncompact}. & Holds for the finitary matched many-sorted presentation by first-order compactness. \\
\addlinespace
Downward L\"owenheim--Skolem & Fails in general; full categoricity of the uncountable ordered field obstructs countable models. & The first-order cardinal-qualified form transfers through the many-sorted presentation. \\
\addlinespace
Upward L\"owenheim--Skolem & Fails in general; full categoricity of \(\mathsf{PA}_2\) obstructs uncountable models. & The first-order form transfers when the object sort \(D\), equivalently the object domain \(M\), is infinite, with the usual language/cardinality qualifications. \\
\addlinespace
Categoricity & Can characterize infinite structures such as \(\mathbb N\) and \(\mathbb R\) up to isomorphism. & The same written axioms need not be categorical; nonstandard selected-domain models remain. \\
\bottomrule
\end{longtable}
\endgroup
\end{landscape}

In the Henkin completeness row, ``complete'' means complete for the exact
calculus/model-class pair; it never means complete for full second-order
validity.  The comparison synthesizes
\cite{Henkin1950,Manzano1996,Shapiro1991}.

\section{Simple type theory and higher-order semantics}
\label{sec:p02-simple-type-theory}

The phrase \emph{higher-order logic} does not determine a calculus.  One must
specify types, term constructors, equality, logical constants, conversion
rules, axioms, and semantics.  The representative language below fixes syntax
and two semantic classes but no local proof calculus.

\begin{hbdefinition}{def:p02-simple-types}{Simple types and typed terms}
Let \(\iota\) be the type of individuals and \(o\) the type of truth values.
The simple types are freely generated by
\[
 \alpha,\beta ::= \iota\mid o\mid(\alpha\mathbin{\to_{\mathrm{ty}}}\beta).
\]
Thus the base types are distinct, arrow formation is injective, and no finite
simple type equals its own function-to-truth type.  For every type, fix typed
variables and possibly typed constants.  If
\(F:\alpha\mathbin{\to_{\mathrm{ty}}}\beta\) and \(A:\alpha\), then
\(FA:\beta\); if \(x:\alpha\) and \(B:\beta\), then
\(\lambda x{:}\alpha.B:\alpha\mathbin{\to_{\mathrm{ty}}}\beta\).
Formulas are exactly terms of type \(o\).  Equality and universal
quantification at type \(\alpha\) have types
\[
 \mathsf{Eq}_{\alpha}:\alpha\mathbin{\to_{\mathrm{ty}}}\alpha
             \mathbin{\to_{\mathrm{ty}}}o,
 \qquad
 \forall_{\alpha}:(\alpha\mathbin{\to_{\mathrm{ty}}}o)
              \mathbin{\to_{\mathrm{ty}}}o.
\]
Thus \(\forall x{:}\alpha.\,A\) abbreviates
\(\forall_{\alpha}(\lambda x{:}\alpha.A)\).
Here \(\alpha,\beta,\gamma\) are metalinguistic type variables, not quantified
object-language type variables, while \(A,B,C,F\) are metalinguistic typed-term
metavariables whenever a type is displayed.  The object-language constant
\(\mathsf{Eq}_{\alpha}\) is distinct from the global metalinguistic
alpha-equivalence relation \(=_{\alpha}\).  The type former
\(\to_{\mathrm{ty}}\) is not object-language implication.  Type arrows
associate to the right and application associates to the left.
\index{simple type theory}
\index{type!simple}
\index{lambda abstraction!typed}
\end{hbdefinition}

Alpha-equivalence, capture-avoiding substitution, beta reduction, and eta
conversion are typed operations.  A mismatched application is not a false
formula; it is not a term.

\begin{hbdefinition}{def:p02-type-structures}{Standard/Henkin type structures}
A \emph{standard type structure} assigns a nonempty set \(D_\alpha\) to each
type, with
\[
 D_o=\{0,1\},
 \qquad
 D_{\alpha\mathbin{\to_{\mathrm{ty}}}\beta}
   =D_\beta^{D_\alpha}.
\]
A \emph{Henkin type structure} instead chooses a nonempty
\[
 D_{\alpha\mathbin{\to_{\mathrm{ty}}}\beta}
 \subseteq D_\beta^{D_\alpha}
\]
that contains the denotation of every well-typed lambda abstraction under
every matching assignment.  In both cases application is function
application, equality is identity, logical constants have their classical
meanings, and a typed assignment sends \(x:\alpha\) to \(D_\alpha\).
Every admitted nonlogical constant \(c:\alpha\) is assigned an element
\(c^{\mathcal M}\in D_\alpha\).
Every standard type structure is a Henkin type structure; the converse fails
when a selected function domain omits set-theoretic functions.
\index{standard semantics!simple type theory}
\index{Henkin semantics!simple type theory}
\end{hbdefinition}

Denotation is the typed recursion
\[
\begin{aligned}
 \lbrack\!\lbrack x\rbrack\!\rbrack_s^{\mathcal M}&=s(x),
&\lbrack\!\lbrack c\rbrack\!\rbrack_s^{\mathcal M}&=c^{\mathcal M},\\
 \lbrack\!\lbrack FA\rbrack\!\rbrack_s^{\mathcal M}
 &=\lbrack\!\lbrack F\rbrack\!\rbrack_s^{\mathcal M}
   (\lbrack\!\lbrack A\rbrack\!\rbrack_s^{\mathcal M}),
&\lbrack\!\lbrack\lambda x{:}\alpha.B\rbrack\!\rbrack_s^{\mathcal M}
 &=\bigl(a\mapsto\lbrack\!\lbrack B\rbrack\!\rbrack_{s[x\mapsto a]}^{\mathcal M}\bigr).
\end{aligned}
\]
The closure clause in a Henkin type structure is therefore a definedness
condition for lambda abstraction, not an optional completeness assumption.
At every type, \(\mathsf{Eq}_{\alpha}\) denotes the characteristic function of identity
on \(D_\alpha\); and \(\forall_\alpha(P)=1\) exactly when
\(P(a)=1\) for every \(a\in D_\alpha\), for each admitted
\(P\in D_{\alpha\mathbin{\to_{\mathrm{ty}}}o}\).

\begin{hbtheorem}{thm:p02-typed-henkin-completeness}{cited=Henkin1950}{Henkin completeness for typed logic}
For the typed formula system, axioms, rules, conversion principles, and general
structures fixed by Henkin, a formula is a theorem iff it is valid in every
one of those general structures.  This theorem neither makes full-function
semantics complete nor supplies a calculus for the syntax-only presentation
of Definition~\ref{def:p02-simple-types}.
\end{hbtheorem}

\begin{hblemma}{thm:p02-typed-substitution-beta}{sketched=annex:p2.hol.9}{Typed substitution and beta soundness}
Let \(B:\beta\), \(A:\alpha\), and \(x:\alpha\), and let \(\mathcal M\)
be a standard or Henkin type structure.  For every matching assignment \(s\),
\[
 \lbrack\!\lbrack B[A/x]\rbrack\!\rbrack_s^{\mathcal M}
 =\lbrack\!\lbrack B\rbrack\!\rbrack_{s[x\mapsto
      \lbrack\!\lbrack A\rbrack\!\rbrack_s^{\mathcal M}]}^{\mathcal M}.
\]
Consequently the beta redex \((\lambda x{:}\alpha.B)A\) and its contractum
\(B[A/x]\) have the same denotation in either semantic class.
\end{hblemma}

Quantification exposes the difference: a standard structure interprets
\(\forall f{:}(\alpha\mathbin{\to_{\mathrm{ty}}}\beta).C\) over all of
\(D_\beta^{D_\alpha}\), while a Henkin structure uses only its declared
\(D_{\alpha\mathbin{\to_{\mathrm{ty}}}\beta}\).  Beta soundness survives
because the latter remains closed under denotations of typed abstractions.

\begin{proofboundary}{def:p02-typed-completeness-boundary}{Typed completeness needs a matched system}
An internal completeness theorem would have to list every primitive typed
constant, conversion rule, logical axiom, equality and extensionality
principle, inference rule, structural convention, and semantic closure
condition, then prove a Henkin construction and truth lemma for exactly that
package.  This chapter does not reproduce those data, introduces no local
typed turnstile, and confines its cited completeness claim to Henkin's source
system \cite{Henkin1950,Church1940}.
\end{proofboundary}

\begin{hbexample}{def:p02-simple-types-polymorphism}{Metalinguistic type schema}
For each fixed simple type \(\alpha\), the identity term
\(\lambda x{:}\alpha.x\) has type
\(\alpha\mathbin{\to_{\mathrm{ty}}}\alpha\).  This is a metalinguistic
schema indexed by \(\alpha\), not one object-language term quantifying over
types.  Object-language type quantification requires a different polymorphic
system.
\end{hbexample}

\begin{hbexample}{def:p02-simple-types-relation}{A typed binary relation}
A binary relation on individuals has type
\(\iota\mathbin{\to_{\mathrm{ty}}}\iota
\mathbin{\to_{\mathrm{ty}}}o\).  If \(R\) has that type and
\(a,b:\iota\), then \(Rab:o\) is a formula; \(Ra\) and \(R\) are not
formulas because their types still end in an unapplied function type.
\end{hbexample}

\begin{hbwarning}{def:p02-type-theory-scope}{Type discipline is not a blanket consistency theorem}
Simple types block unrestricted self-application and thereby the syntactic
form of Russell's self-membership paradox.  They do not prove consistency of
every extension.  Functional extensionality is the principle
\[
 \forall f,g{:}(\alpha\mathbin{\to_{\mathrm{ty}}}\beta).\,
 \bigl(\forall x{:}\alpha.\,\mathsf{Eq}_{\beta}\,(fx)\,(gx)\bigr)\to
 \mathsf{Eq}_{\alpha\mathbin{\to_{\mathrm{ty}}}\beta}\,f\,g.
\]
A calculus may take it as an axiom, while a semantics may validate it by
interpreting equality as identity of functions.  Functional extensionality
and higher-type choice are separate principles, not consequences of the type
grammar.  Likewise, a proof assistant's metalogic proves nothing about an
encoded object logic until a theorem connects encoded derivability or
satisfaction to the intended external judgment.
\end{hbwarning}
